\uuid{VZFM}
\titre{Une fonction de répartition donnée}
\auteur{Maxime Nguyen}
\organisation{AMSCC}
\datecreate{2026-11-24}
\niveau{L2}
\module{Probabilité}
\chapitre{Probabilité continue}
\sousChapitre{Espérance et Fonction de répartition}
\difficulte{2}
\theme{fonction de répartition, densité, espérance, variance, médiane}
\serie{ordonnee}

\contenu{
\texte{
On considère la fonction $F$ définie par
$$F(t) = \begin{cases} 0 & \text{si } t < 1,\\[2pt] 1 - \dfrac{1}{t^2} & \text{si } t \geq 1.\end{cases}$$
On admet que $F$ est la fonction de répartition d'une variable à densité $X$. Dans les réponses, $\mathbf{1}_{[1;+\infty[}$ désigne l'indicatrice de $[1;+\infty[$ : elle vaut $1$ sur cet intervalle et $0$ ailleurs.
}

\begin{question}
\qtype{onlyone}
Une densité de $X$ est :
\begin{answers}
\good{$f(x) = \dfrac{2}{x^3}\,\mathbf{1}_{[1;+\infty[}(x)$}
\bad{$f(x) = \dfrac{2}{x^3}$ pour tout $x \neq 0$}
\bad{$f(x) = \left(1 - \dfrac{1}{x^2}\right)\mathbf{1}_{[1;+\infty[}(x)$}
\bad{$X$ n'a pas de densité, car $F$ n'est pas dérivable en $1$.}
\end{answers}
\begin{explanations}
$F'(t) = \frac{2}{t^3}$ pour $t > 1$ et $F'(t) = 0$ pour $t < 1$.

B oublie le support : on a dérivé la formule $1 - \frac{1}{t^2}$ hors de son
domaine de validité, et la fonction obtenue est négative pour $x < 0$ et non
intégrable près de $0$. C est la fonction de répartition elle-même. D : $F$ est
continue et $\mathcal{C}^1$ par morceaux, un point de non-dérivabilité
n'empêche pas l'existence d'une densité.

Contre-sens visé : dériver une formule sans tenir compte de son domaine.
\end{explanations}
\end{question}

\begin{question}
\qtype{onlyone}
$\mathbb{P}(X > 3)$ vaut :
\begin{answers}
\bad{$\frac{8}{9}$}
\bad{$\frac{25}{27}$}
\bad{$\frac{2}{27}$}
\good{$\frac{1}{9}$}
\end{answers}
\begin{explanations}
$\mathbb{P}(X > 3) = 1 - F(3) = \frac19$.

A donne $F(3) = \mathbb{P}(X \leq 3)$ ; C donne $f(3)$ et B donne $1 - f(3)$,
deux calculs qui prennent la densité pour une probabilité.

Contre-sens visé : confondre $F(t)$, $1 - F(t)$ et $f(t)$.
\end{explanations}
\end{question}

\begin{question}
\qtype{onlyone}
$\mathbb{E}(X)$ vaut :
\begin{answers}
\good{$2$}
\bad{$1$}
\bad{$\sqrt2$}
\bad{$+\infty$ : $X$ n'a pas d'espérance.}
\end{answers}
\begin{explanations}
$\mathbb{E}(X) = \int_1^{+\infty} x\cdot\frac{2}{x^3}\,\mathrm{d}x = \int_1^{+\infty}\frac{2}{x^2}\,\mathrm{d}x = 2$.

B correspond à $\int f = 1$ (oubli du facteur $x$). C est la médiane :
$F(t) = \frac12$ pour $t = \sqrt2$. D : la queue de la densité en
$\frac{2}{x^3}$ décroît assez vite pour que $x f(x)$ soit intégrable.

Contre-sens visé : confondre espérance et médiane.
\end{explanations}
\end{question}

\begin{question}
Quelles affirmations sont vraies ?
\begin{answers}
\bad{$X$ admet une variance, que l'on calcule par $\mathbb{E}(X^2) - \mathbb{E}(X)^2$.}
\good{$\mathbb{E}(X^2) = +\infty$, donc $X$ n'admet pas de variance.}
\bad{Une variable qui admet une espérance admet toujours une variance.}
\end{answers}
\begin{explanations}
$\int_1^{+\infty} x^2\cdot\frac{2}{x^3}\,\mathrm{d}x = \int_1^{+\infty}\frac{2}{x}\,\mathrm{d}x = +\infty$.
C est faux : cette variable en est un contre-exemple.

Contre-sens visé : croire que toute variable a une espérance et une variance finies.
\end{explanations}
\end{question}
}
